\section{Introduction}

Human-object interaction (HOI) animation, which involves generating dynamic motion sequences of a person interacting with objects, has become an important problem in computer vision and graphics. Realistic HOI animation underpins applications ranging from virtual reality, gaming, and robotics simulation, where agents must manipulate or use objects in a human-like way. Recent advances in generative model have opened new possibilities for synthesizing such animations directly from high-level specifications, \textit{e.g.,} text descriptions. In particular, diffusion models~\citep{peng2023hoi,diller2023cg,li2023controllable} have emerged as a powerful generative framework capable of modeling high-dimensional distribution of human motions and contacts, by iteratively denoising noise into plausible HOI sequences.

Despite this promise, a persistent challenge is ensuring high-quality realistic interactions between humans and objects. A plain diffusion model, without physics simulation, often produces artifacts: hands miss their targets, objects drift or penetrate the body, or contact is unstable over time. Prior work has sought to mitigate this through \textit{guidance}. One line of methods employs external classifiers trained for \textit{human-specified tasks}, \textit{e.g.}, contact or affordance regression~\citep{peng2023hoi,diller2023cg} to steer the diffusion process, but such classifier guidance is cumbersome to design, and may overfit to specific priors. Another line introduces \textit{hand-crafted rules} with kinematic and dynamic constraints, such as forcing hands to align with objects via inverse kinematics~\citep{xu2024interdreamer}, or introducing physics simulation~\citep{xu2025intermimic,xu2026interprior}, but these approaches sacrifice generality and are computationally inefficient. Together, these efforts highlight: existing HOI animation still relies heavily on external priors not directly from the data.

Extending classifier-free guidance (CFG)~\citep{ho2022classifier} to HOI animation is a natural way to reduce the reliance on external priors, but CFG, especially CFG based on text dropout for text-conditioned generation, mainly improves global distributional alignment and offers limited control over the fine-grained, persistent contact central to HOI.
We therefore introduce LIGHT, \textit{i.e.,} \textit{L}earning \textit{I}mplicit \textit{G}uidance for \textit{H}uman-object in\textit{T}eraction, a complementary data-driven framework where guidance arises from the relative denoising pace of separate components. As illustrated in Figure~\ref{fig:teaser}, LIGHT produces diverse and realistic human-object interaction animations with coherent contact dynamics. Concretely, we factor the representation into modality-specific components (\textit{e.g.,} human and object) and assign each its own noise level with asynchronous schedules. And we formulate two paths:
(\textbf{I}) a \textit{staged schedule}, where one modality (\textit{e.g.}, human) is kept cleaner while the other (\textit{e.g.}, object) follows its prescribed noise levels, approximating a conditional trajectory; and
(\textbf{II}) a \textit{uniform schedule} that applies weak conditioning by assigning all modalities with prescribed noise levels, where the ``unconditional'' case is realized as noisier conditioning, constituting a \textit{soft} form of CFG. The contrast between these two paths produces a guidance effect analogous to CFG. As the lag between schedules approaches zero, LIGHT collapses to joint denoising with no guidance; as the lag grows large, it approximates conditioning dropout in CFG. Our experiment shows that hard dropout on text improves global distributional alignment, whereas the soft guidance from LIGHT adjusts more on low-level contact details, indicating that LIGHT learns, purely from data, to reduce contact errors without hand-crafted priors. Note that assigning different noise levels to diffusion models naturally corresponds to the diffusion forcing mechanism~\citep{chen2024diffusion}, and LIGHT provides a principled extension of diffusion forcing into a guidance framework.

LIGHT relies solely on data priors; however, this does not constrain the model’s flexibility to incorporate additional world knowledge like physics provided by humans, as such human-induced priors can also be reflected through data manipulation. We find that LIGHT achieves further improvement when training data are augmented with prior via \textit{contact-aware shape-spectrum augmentation}. Specifically, objects are augmented with geometrically diverse alternatives from large repositories~\citep{chang2015shapenet, deitke2023objaverse}, and motions are retargeted to preserve contact semantics and relative pose. These synthetic pairs teach the model that contact should remain invariant to irrelevant shape changes. This results in a stronger prior that the asynchronous schedules can exploit, improving generalizability to, \textit{e.g.,}  unseen objects from training.

In summary, our contribution lies in three folds:
\begin{itemize}
    \item We introduce LIGHT, a novel guidance mechanism in which asynchronous denoising induces guidance without reliance on external classifiers. In contrast to CFG’s reliance on hard dropout, our formulation provides a softer and more flexible alternative, naturally extending diffusion forcing into the guidance framework and potentially \textit{inspiring future applications} in other domains.
    \item We propose contact-aware shape-spectrum augmentation, a strategy that preserves contact semantics while varying object geometry, thereby enabling the model to acquire more robust generative capabilities and improve generalization.
    \item We conduct extensive experimental evaluations demonstrating that the proposed approach consistently outperforms existing baselines and facilitates the generation of vivid and realistic interactions.
\end{itemize}